\section{Geometric two-view baselines}
\label{sec:geometric}

All three reviews ask how correspondence-based pipelines perform on this task: local-feature matchers followed by robust two-view geometry, feed-forward reconstruction models used on their own, and geometry combined with a metric scale. The main paper compares with none of them. This section adds them on PairUAV; \cref{sec:real} adds them on the real photographs. The outcome changes how the results of the main paper should be read, and we state it plainly: \emph{a classical SIFT pipeline with the same flight-height prior is more accurate than our coarse stage and within about five points of our refined stage on PairUAV; a recent dense matcher is level with the refined stage on the pairs on which we could run it; and both are more accurate than DreamNav on the real photographs.}

\subsection{Protocol}
\label{sec:geoprotocol}

\noindent\textbf{Methods.} All methods use public weights; nothing is trained on PairUAV. We evaluate SIFT with the ratio test~\cite{lowe2004sift}; SuperPoint~\cite{detone2018superpoint} with SuperGlue~\cite{sarlin2020superglue} and with LightGlue~\cite{lindenberger2023lightglue}; ALIKED~\cite{zhao2023aliked} with LightGlue; LoFTR~\cite{sun2021loftr}; DeDoDe~\cite{edstedt2024dedode} with a rotation steerer~\cite{bokman2024steerers}; RoMa~\cite{edstedt2024roma}; the dense descriptors of MASt3R~\cite{leroy2024mast3r} with reciprocal matching; and the pointmaps of MASt3R and DUSt3R~\cite{wang2024dust3r} followed by PnP. The learned matchers run at the resolution of their public configuration (long side of 1{,}024 to 1{,}152 pixels, up to 2{,}048 keypoints for the sparse ones); SIFT runs on the half-resolution image (1{,}843 pixels) with up to 8{,}000 features, a ratio of 0.8 and mutual nearest neighbours. RoMa and the MASt3R descriptors are slow on the shared GPU and were run on a random subset of the pairs, on which our own rows are also given.

\noindent\textbf{Pose.} Every matcher feeds the same solver, which is the one of our coarse stage without its priors, its refinement and its fusion (\cref{sec:coarse}): a calibrated essential matrix estimated with MAGSAC++~\cite{barath2020magsacpp}, the choice among its four decompositions by cheirality and by the zero-roll, limited-tilt convention of the dataset, and the decomposition into $\dth$, $\dpsi$ and a unit translation. The inlier threshold is selected for each matcher among seven values on \gbNcal{} training pairs. A pair for which no pose is found predicts zero motion.

\noindent\textbf{In-plane rotation.} The two views of a pair have arbitrary headings. SIFT is invariant to this rotation and DeDoDe with the steerer is equivariant. For the other matchers we report the result on the pair as it is and with a test-time search: the target is matched under its four rotations by multiples of $90^\circ$ and the rotation with the most inliers is kept, which is what the rotation search of our first pass does with its own criterion.

\noindent\textbf{Metric scale.} An essential matrix leaves the length of the translation undetermined. Every geometric row takes it from the flight-height prior that our system uses as one of its cues (\cref{sec:cands}): the height of the two cameras above the triangulated ground, in units of the baseline, is equated with a constant $H$ that is fitted on the same \gbNcal{} training pairs (218.9\,m for SIFT). Other sources of scale are compared in \cref{tab:geo_scale}.

\begin{table*}[t]
\centering
\caption{\textbf{Geometric two-view baselines on PairUAV} (1{,}000 random test pairs, a random subset of them where fewer are listed; inlier thresholds and the constant of the flight-height prior from 300 training pairs). Camera frame. Rot.: handling of the in-plane rotation (inv.: invariant descriptor; eq.: rotation-equivariant matching; TTA: the target is matched under four rotations and the one with most inliers is kept; ours: rotation search of Sec.~4.2). Fail: pairs without a pose, without and with this handling; they predict zero motion. Angular errors are mean / median.}
\label{tab:geometric}
\footnotesize
\setlength{\tabcolsep}{3.5pt}
\resizebox{\textwidth}{!}{%
\begin{tabular}{@{}llrcccccccccc@{}}
\toprule
& & & \multicolumn{2}{c}{Fail (\%)} & $d$ MAE & \SRd@1\,m & \multicolumn{2}{c}{EPE (m)} & \SRe@1\,m & $\dth$ error ($^\circ$) & $\dpsi$ error ($^\circ$) & Dir. \\
Method & Rot. & Pairs & as is & handled & (m) & (\%) & mean & med. & (\%) & mean / med. & mean / med. & med. ($^\circ$) \\
\midrule
\multicolumn{13}{@{}l}{\emph{Matcher + calibrated essential matrix (MAGSAC++); length from the flight-height prior}}\\
SIFT + ratio test & inv. & 1{,}000 & -- & 0.0 & 0.89 & 78.9 \ci{76.1}{81.4} & 1.13 & 0.61 & 74.2 \ci{71.1}{76.9} & 0.051 / 0.013 & 0.249 / 0.058 & 0.5 \\
SuperPoint + SuperGlue, one pass & none & 1{,}000 & 55.3 & -- & 12.31 & 24.6 \ci{21.9}{27.1} & 12.95 & 12.06 & 16.1 \ci{13.7}{18.4} & 65.9 / 72.5 & 11.4 / 4.09 & 90.0 \\
SuperPoint + SuperGlue, our first pass & ours & 1{,}000 & 55.3 & 0.7 & 3.09 & 56.7 \ci{53.6}{59.8} & 4.78 & 1.13 & 45.3 \ci{41.9}{48.7} & 1.46 / 0.032 & 1.61 / 0.164 & 1.6 \\
SuperPoint + LightGlue & TTA & 1{,}000 & 61.3 & 1.9 & 2.72 & 51.3 \ci{48.2}{54.4} & 4.28 & 1.56 & 35.3 \ci{32.3}{38.2} & 1.78 / 0.053 & 1.77 / 0.260 & 2.7 \\
ALIKED + LightGlue & TTA & 1{,}000 & 51.8 & 3.3 & 4.22 & 41.1 \ci{38.3}{44.3} & 6.62 & 2.75 & 20.5 \ci{17.6}{23.1} & 3.66 / 0.090 & 2.69 / 0.536 & 5.7 \\
LoFTR & TTA & 1{,}000 & 44.3 & 3.5 & 7.40 & 49.9 \ci{46.6}{53.2} & 10.02 & 1.60 & 35.5 \ci{32.5}{38.6} & 6.13 / 0.057 & 3.61 / 0.286 & 2.9 \\
DeDoDe + rotation steerer & eq.+TTA & 1{,}000 & 0.2 & 0.0 & 1.82 & 65.5 \ci{62.5}{68.7} & 2.92 & 0.92 & 54.1 \ci{50.9}{57.3} & 0.158 / 0.021 & 0.830 / 0.113 & 1.1 \\
RoMa & TTA & 200 & 8.0 & 0.0 & 1.56 & 81.5 \ci{75.4}{87.0} & 2.05 & 0.50 & 79.0 \ci{72.9}{84.5} & 0.100 / 0.013 & 0.238 / 0.053 & 0.5 \\
MASt3R descriptors & TTA & 200 & 0.0 & 0.5 & 2.65 & 49.0 \ci{42.0}{55.4} & 4.20 & 2.16 & 23.5 \ci{18.0}{28.9} & 0.235 / 0.076 & 1.37 / 0.481 & 4.3 \\
\midrule\multicolumn{13}{@{}l}{\emph{Feed-forward reconstruction used as a pose estimator (pointmap + PnP)}}\\
MASt3R, its own metric scale & TTA & 1{,}000 & 0.1 & 0.0 & 20.15 & 0.1 \ci{0.0}{0.3} & 20.38 & 20.47 & 0.0 \ci{0.0}{0.0} & 0.654 / 0.375 & 4.90 / 2.48 & 39.6 \\
MASt3R, flight-height prior & TTA & 1{,}000 & 0.4 & 0.3 & 16.91 & 6.8 \ci{5.1}{8.6} & 26.68 & 17.35 & 0.1 \ci{0.0}{0.4} & 0.651 / 0.375 & 4.95 / 2.50 & 39.6 \\
DUSt3R, flight-height prior & TTA & 1{,}000 & 0.1 & 0.1 & 16.99 & 6.4 \ci{5.0}{8.0} & 28.48 & 20.06 & 0.0 \ci{0.0}{0.0} & 0.700 / 0.400 & 4.79 / 2.51 & 49.5 \\
\midrule\multicolumn{13}{@{}l}{\emph{Learned}}\\
Regression on pooled MASt3R statistics &  & 1{,}000 & -- & 0.0 & 3.21 & 20.4 \ci{17.9}{23.1} & 6.76 & 6.11 & 1.0 \ci{0.4}{1.7} & 2.01 / 1.13 & 1.86 / 1.30 & 14.6 \\
Ours (Stage I) &  & 1{,}000 & -- & 0.0 & 0.93 & 73.8 \ci{70.8}{76.5} & 1.67 & 0.82 & 60.5 \ci{57.2}{63.8} & 0.101 / 0.020 & 0.397 / 0.096 & 1.0 \\
Ours (Stage II) &  & 1{,}000 & -- & 0.0 & 0.67 & 83.1 \ci{80.5}{85.7} & 1.01 & 0.52 & 78.4 \ci{75.7}{81.2} & 0.049 / 0.006 & 0.195 / 0.026 & 0.3 \\
\bottomrule
\end{tabular}}
\end{table*}

\begin{table}[t]
\centering
\caption{\textbf{Where the metric scale comes from.} Upper part: translation direction and rotation from SIFT + essential matrix on the pairs of \cref{tab:geometric}; only the source of the length changes. Lower part: our own direction with the true length, for reference.}
\label{tab:geo_scale}
\footnotesize
\setlength{\tabcolsep}{4pt}
\resizebox{\columnwidth}{!}{%
\begin{tabular}{@{}lcccc@{}}
\toprule
& $d$ MAE & \SRd & EPE med. & \SRe \\
Source of the length & (m) & (\%) & (m) & (\%) \\
\midrule
MASt3R metric depth, as is & 20.37 & 0.0 & 20.38 & 0.0 \\
MASt3R metric depth $\times$ one constant & 2.03 & 35.3 & 1.54 & 32.6 \\
Flight-height prior & 0.89 & 78.9 & 0.61 & 74.2 \\
Learned distance regression & 2.69 & 24.7 & 2.20 & 22.7 \\
Flight height and regression, blended & 0.85 & 79.0 & 0.62 & 74.2 \\
\midrule
Fused distance of our Stage I & 0.95 & 73.7 & 0.69 & 68.3 \\
Fused distance of our Stage II & 0.68 & 82.1 & 0.58 & 76.6 \\
True distance (oracle) & 0.00 & 100.0 & 0.19 & 93.4 \\
\midrule\multicolumn{5}{@{}l}{\emph{Direction and rotation of our stages, true distance (oracle)}}\\
Stage I & 0.00 & 100.0 & 0.35 & 78.5 \\
Stage II & 0.00 & 100.0 & 0.12 & 92.1 \\
\bottomrule
\end{tabular}}
\end{table}


\subsection{Results on PairUAV}
\label{sec:georesults}

\noindent\textbf{Rotation handling is a prerequisite.} As they are, SuperGlue, LightGlue, ALIKED and LoFTR return no pose for \gbPlainFailLo\% to \gbPlainFailHi\% of the pairs (\cref{tab:geometric}, ``as is''). With the four-rotation search the failures drop to a few percent. A baseline that is not given this search is not a meaningful reference on this dataset.

\noindent\textbf{A classical pipeline between our two stages.} SIFT, the essential matrix and the flight-height prior reach \gbSiftSRd\% \SRd{} and \gbSiftSRe\% \SRe{} on the pairs of \cref{tab:geometric}, with a heading error of $\gbSiftTheta^\circ$. We therefore ran this pipeline on all \sfNfull{} test pairs of Table~2: it obtains \sfSiftSRd\% \ci{\sfSiftSRdLo}{\sfSiftSRdHi} / \sfSiftSRe\% \ci{\sfSiftSReLo}{\sfSiftSReHi}, where our coarse stage has \sfSiSRd\% / \sfSiSRe\% and our refined stage \sfSiiSRd\% / \sfSiiSRe\%. Paired on the same pairs, it is \sfVsISRd{} \ci{\sfVsISRdLo}{\sfVsISRdHi} and \sfVsISRe{} \ci{\sfVsISReLo}{\sfVsISReHi} points above the coarse stage, and \sfVsIISRd{} \ci{\sfVsIISRdLo}{\sfVsIISRdHi} and \sfVsIISRe{} \ci{\sfVsIISReLo}{\sfVsIISReHi} points from the refined stage. Rotated into ENU with the heading network of the paper, as in Table~2, its endpoint success is \sfSiftEnuSRe\% against \sfSiiEnuSRe\%. The comparison depends on the baseline length: for the \sfShortN{} pairs with less than 10\,m between the cameras the SIFT pipeline is ahead of the refined stage (\sfShortDelta{} points of \SRe{} for ours, \ci{\sfShortDeltaLo}{\sfShortDeltaHi}), and for the \sfLongN{} pairs beyond 30\,m the refined stage is ahead (\sfLongDelta{} \ci{\sfLongDeltaLo}{\sfLongDeltaHi}).

\noindent\textbf{A dense matcher level with the refined stage.} RoMa with the rotation search reaches \gbRomaSRd\% \ci{\gbRomaSRdLo}{\gbRomaSRdHi} / \gbRomaSRe\% \ci{\gbRomaSReLo}{\gbRomaSReHi} on \gbRomaN{} pairs, on which the refined stage has \gbSiiOnRomaSRd\% / \gbSiiOnRomaSRe\%; the paired differences, \gbRomaVsIISRd{} \ci{\gbRomaVsIISRdLo}{\gbRomaVsIISRdHi} and \gbRomaVsIISRe{} \ci{\gbRomaVsIISReLo}{\gbRomaVsIISReHi} points, are not significant at this sample size. Without the rotation search it finds no pose for \gbRomaPlainFail\% of the pairs and reaches \gbRomaPlainSRd\%.

\noindent\textbf{Two settings matter for the baselines.} The inlier threshold and the statistic of the flight-height prior were chosen on the training pairs, and the result is sensitive to both. With the 3\,px threshold of our own first pass instead of the selected 0.2\,px, the SIFT pipeline drops to \gbSiftThrThreeSRd\% \SRd; with the median of the triangulated heights instead of their 75th percentile, which is closer to the ground, it drops to \gbSiftMedHeightSRd\%. A geometric baseline assembled without this care would have looked much weaker than it is.

\noindent\textbf{Learned matchers.} At their standard resolution the sparse and semi-dense learned matchers are less accurate than SIFT on these images: DeDoDe with the steerer reaches \gbDedodeSRd\% / \gbDedodeSRe\%, LightGlue \gbLgSRd\% / \gbLgSRe\%, LoFTR \gbLoftrSRd\% / \gbLoftrSRe\% and ALIKED \gbAlikedSRd\% / \gbAlikedSRe\%. Their median error on the translation direction is $\gbDedodeDir^\circ$ to $\gbAlikedDir^\circ$, against $\gbSiftDir^\circ$ for SIFT and $\gbSiiDir^\circ$ for our refined stage. The baseline of a pair subtends about $5^\circ$ at the ground, so that a success at 1\,m requires correspondences that are accurate to about a pixel of the original image; SIFT localises its keypoints with sub-pixel accuracy on an image almost twice as large as the input of the learned matchers. We did not tune the learned matchers beyond their public configuration, and they may improve at a higher resolution. The observation agrees with the controls of \cref{sec:controls}, where the resolution of the second matching pass, not its alignment model, carries most of the gain of our refinement.

\noindent\textbf{Feed-forward reconstruction.} Used directly as a two-view pose estimator, MASt3R recovers the rotation (mean heading error $\gbMastTheta^\circ$) but not the translation: the median direction error is $\gbMastDir^\circ$, and its metric scale is wrong by a factor of about one hundred on views taken from 220\,m, so that its error equals that of predicting no motion (\gbMastNativeEPE\,m). With the flight-height prior it reaches \gbMastHeightSRd\% \SRd, and DUSt3R \gbDustHeightSRd\%. Matching the dense MASt3R descriptors and solving the essential matrix is better (\gbMastMatchSRd\% / \gbMastMatchSRe\% on \gbMastMatchN{} pairs) but remains far from the keypoint-based pipelines. The models are evaluated at their 512-pixel input, at which one metre of displacement is about two pixels.

\noindent\textbf{The scale comes from the flight height, for every method.} \Cref{tab:geo_scale} keeps the direction and rotation of the SIFT pipeline and changes the source of the length. The metric depth of MASt3R gives \gbScaleNativeSRd\% as it is and \gbScaleCalibSRd\% after multiplication by one constant; a learned regression of the distance gives \gbScaleLearnedSRd\%; the flight-height prior gives \gbScaleHeightSRd\%. The geometric baselines are therefore metric only because the flight height of PairUAV is regular (220\,m\,$\pm$\,9\,m, \cref{sec:dataset}), which is also the main source of scale of our own system (\cref{sec:inventory}).

\noindent\textbf{Where the remaining difference lies.} With the true length, the direction and rotation of the SIFT pipeline succeed on \gbHybOracleSRe\% of the pairs, those of our refined stage on \gbOracleSiiSRe\% and those of our coarse stage on \gbOracleSiSRe\%. The refinement therefore brings our direction to the level of the classical pipeline, not beyond it. The same holds without any learning: the match set of our refined stage, given to the same solver with the same flight-height prior, reaches \gFiveSRd\% / \gFiveSRe\% on the full test set (row G5 of \cref{tab:inventory}), against \sfSiftSRd\% / \sfSiftSRe\% for SIFT. With the fused distance of our refined stage in place of its own flight-height estimate, the SIFT pipeline reaches \gbHybTwoSRd\% / \gbHybTwoSRe\%. The advantage of the refined stage over the classical pipeline on PairUAV is thus mostly an advantage in the estimated \emph{distance} (range error \gbHybTwoD\,m against \gbSiftD\,m), obtained by fusing several scale cues with a learned regressor. The two systems also fail on different pairs: on the full test set both succeed (endpoint error below 1\,m) on \sfJointBoth\% of the pairs, only the SIFT pipeline on \sfJointSift\%, only ours on \sfJointOurs\%, neither on \sfJointNone\%, so that the better of the two would reach \sfJointEither\%.

\begin{table}[t]
\centering
\caption{\textbf{Pairs that were not screened.} 1{,}160 ordered pairs drawn uniformly from all camera pairs of the 290 test buildings. Camera frame. Fail: pairs without a pose (zero motion is predicted).}
\label{tab:unscreened}
\footnotesize
\setlength{\tabcolsep}{3pt}
\resizebox{\columnwidth}{!}{%
\begin{tabular}{@{}lccccc@{}}
\toprule
& Fail & \SRd@1\,m & EPE med. & \SRe@1\,m & $\dth$ med. \\
& (\%) & (\%) & (m) & (\%) & ($^\circ$) \\
\midrule
\multicolumn{6}{@{}l}{\emph{All sampled pairs} (1{,}160 pairs)}\\
\quad SIFT + E, flight height & 19.5 & 55.0 \ci{52.1}{58.3} & 0.92 & 52.6 \ci{49.5}{56.0} & 0.021 \\
\quad Ours (Stage I) & 0.0 & 54.1 \ci{51.0}{57.1} & 1.16 & 45.9 \ci{42.9}{48.9} & 0.028 \\
\quad Ours (Stage II) & 0.0 & 60.9 \ci{58.1}{64.0} & 0.84 & 55.4 \ci{52.3}{58.7} & 0.010 \\
\midrule\multicolumn{6}{@{}l}{\emph{Within the released ranges, passing the screen} (490 pairs)}\\
\quad SIFT + E, flight height & 0.2 & 83.5 \ci{80.0}{86.6} & 0.50 & 79.8 \ci{76.3}{83.0} & 0.013 \\
\quad Ours (Stage I) & 0.0 & 80.0 \ci{76.4}{83.5} & 0.64 & 72.4 \ci{68.7}{76.5} & 0.012 \\
\quad Ours (Stage II) & 0.0 & 86.3 \ci{83.2}{89.4} & 0.49 & 82.2 \ci{78.8}{85.6} & 0.005 \\
\midrule\multicolumn{6}{@{}l}{\emph{Within the released ranges} (704 pairs)}\\
\quad SIFT + E, flight height & 18.2 & 61.4 \ci{57.5}{65.0} & 0.75 & 58.1 \ci{54.2}{62.0} & 0.020 \\
\quad Ours (Stage I) & 0.0 & 62.4 \ci{58.6}{66.0} & 0.95 & 52.0 \ci{48.4}{55.7} & 0.027 \\
\quad Ours (Stage II) & 0.0 & 67.6 \ci{64.0}{71.1} & 0.71 & 60.8 \ci{57.2}{64.4} & 0.010 \\
\midrule\multicolumn{6}{@{}l}{\emph{Outside the released ranges} (456 pairs)}\\
\quad SIFT + E, flight height & 21.5 & 45.2 \ci{40.6}{50.0} & 1.13 & 44.1 \ci{39.6}{48.7} & 0.023 \\
\quad Ours (Stage I) & 0.0 & 41.2 \ci{37.0}{45.7} & 1.51 & 36.4 \ci{32.1}{40.7} & 0.028 \\
\quad Ours (Stage II) & 0.0 & 50.7 \ci{46.2}{55.1} & 1.13 & 47.1 \ci{42.5}{51.8} & 0.009 \\
\midrule\multicolumn{6}{@{}l}{\emph{Failing the screen} (366 pairs)}\\
\quad SIFT + E, flight height & 61.5 & 10.7 \ci{8.0}{13.6} & 17.22 & 8.2 \ci{5.5}{11.0} & 125.4 \\
\quad Ours (Stage I) & 0.0 & 16.1 \ci{12.5}{19.8} & 18.56 & 3.8 \ci{2.0}{5.9} & 4.93 \\
\quad Ours (Stage II) & 0.0 & 18.9 \ci{14.6}{22.7} & 17.32 & 9.3 \ci{6.5}{12.1} & 4.72 \\
\bottomrule
\end{tabular}}
\end{table}


\subsection{Is SIFT favoured by the screening of the dataset?}
\label{sec:unscreened}

Every released pair passed an overlap screen that is itself based on SIFT matches (\cref{sec:dataset}), so that SIFT matching succeeds on all test pairs by construction. To measure how much this helps the baseline, we draw \unN{} ordered pairs uniformly from all camera pairs of the test buildings, four per building, without any screening. Of these pairs, \unInside\% lie within the ranges of pose difference of the released pairs, \unScreenPass\% pass our re-implementation of the screen, and \unNoOverlap\% have no common ground footprint at all.

\Cref{tab:unscreened} gives the result. The SIFT pipeline returns no pose for \unSiftFail\% of the pairs, whereas ours always returns an estimate; over all pairs the refined stage succeeds on \unSiiSRd\% / \unSiiSRe\% and the SIFT pipeline on \unSiftSRd\% / \unSiftSRe\%, a paired difference of \unSiiVsSiftSRd{} \ci{\unSiiVsSiftSRdLo}{\unSiiVsSiftSRdHi} and \unSiiVsSiftSRe{} \ci{\unSiiVsSiftSReLo}{\unSiiVsSiftSReHi} points. The margin is the same as on the released pairs, so the screen does not explain the strength of the baseline. On the pairs that fail the screen, which mostly share little or no ground, neither system is usable (\unFailSiiSRe\% and \unFailSiftSRe\% \SRe): an estimate that is always returned is not, on these pairs, an estimate that can be trusted. The table also quantifies the effect of the screening on our own numbers, which the reviews ask for: without it, the success of the refined stage is \unSiiSRd\% instead of \SiiCamSRd\%.

\begin{table}[t]
\centering
\caption{\textbf{The SIFT pose and scale cues as additional candidates of our fusion} (3{,}000 test pairs; fusions fitted on 3{,}000 training pairs, as in \cref{tab:controls}). Camera frame.}
\label{tab:sift_fusion}
\footnotesize
\setlength{\tabcolsep}{3pt}
\resizebox{\columnwidth}{!}{%
\begin{tabular}{@{}lccccc@{}}
\toprule
& $d$ & \SRd@1\,m & EPE med. & \SRe@1\,m & Dir. med. \\
& (m) & (\%) & (m) & (\%) & ($^\circ$) \\
\midrule
SIFT + E, flight-height prior & 1.06 & 78.4 \ci{76.3}{80.1} & 0.61 & 73.7 \ci{71.7}{75.6} & 0.56 \\
\midrule Ours, coarse stage & 0.91 & 73.5 \ci{71.8}{75.1} & 0.82 & 60.5 \ci{58.3}{62.6} & 1.01 \\
\quad + SIFT candidates & 0.70 & 81.8 \ci{80.1}{83.3} & 0.59 & 76.6 \ci{74.6}{78.3} & 0.58 \\
\midrule Ours, refined stage & 0.70 & 82.7 \ci{80.9}{84.4} & 0.54 & 78.3 \ci{76.2}{80.1} & 0.41 \\
\quad + SIFT candidates & 0.64 & 84.1 \ci{82.5}{85.6} & 0.51 & 80.6 \ci{78.7}{82.2} & 0.38 \\
\bottomrule
\end{tabular}}
\end{table}


\subsection{Combining the classical pipeline with our fusion}
\label{sec:siftfusion}

Since the two systems fail on different pairs, we add the pose and the scale cues of the SIFT pipeline to the candidates of our fusion (\cref{sec:fusion}) and fit it again, on the \sfuNtrain{} training pairs of the controls of \cref{sec:controls}; \cref{tab:sift_fusion} reports their \sfuN{} test pairs.

\noindent\textbf{SIFT candidates in place of the imagination.} Added to the coarse stage, with no second matching pass of ours, the SIFT candidates raise \SRd{} from \sfuCoarseSRd\% to \sfuCoarseSiftSRd\% and \SRe{} from \sfuCoarseSRe\% to \sfuCoarseSiftSRe\% (\sfuAddCoarseSRd{} \ci{\sfuAddCoarseSRdLo}{\sfuAddCoarseSRdHi} and \sfuAddCoarseSRe{} \ci{\sfuAddCoarseSReLo}{\sfuAddCoarseSReHi} points). The imagination-refined stage is \sfuRefVsCoarseSiftSRd{} \ci{\sfuRefVsCoarseSiftSRdLo}{\sfuRefVsCoarseSiftSRdHi} points of \SRd{} and \sfuRefVsCoarseSiftSRe{} \ci{\sfuRefVsCoarseSiftSReLo}{\sfuRefVsCoarseSiftSReHi} points of \SRe{} above this configuration. An independent set of accurate correspondences from a classical detector thus gives the fusion nearly the information that the imagined view gives it. This is a further control for the question of \cref{sec:controls}: most of what the refinement stage contributes is a second, more accurate match set, however it is obtained.

\noindent\textbf{SIFT candidates in addition to the imagination.} Added to the refined stage, they bring \sfuAddRefSRd{} \ci{\sfuAddRefSRdLo}{\sfuAddRefSRdHi} and \sfuAddRefSRe{} \ci{\sfuAddRefSReLo}{\sfuAddRefSReHi} further points, to \sfuRefSiftSRd\% / \sfuRefSiftSRe\%, the best result that we obtained on PairUAV. The two match sets are therefore partly complementary.

\noindent\textbf{What the fusion adds to SIFT.} With the same SIFT matches, the fusion with the coarse stage is \sfuCoarseSiftVsSiftSRd{} \ci{\sfuCoarseSiftVsSiftSRdLo}{\sfuCoarseSiftVsSiftSRdHi} / \sfuCoarseSiftVsSiftSRe{} \ci{\sfuCoarseSiftVsSiftSReLo}{\sfuCoarseSiftVsSiftSReHi} points above the SIFT pipeline alone, and the complete system \sfuAllVsSiftSRd{} / \sfuAllVsSiftSRe{} points. The learned combination of several pose and scale estimates is the component of DreamNav that improves on classical geometry on PairUAV; we have not tested whether this holds on the real photographs, where the fusion would have to be re-fitted.

\subsection{Consequences}
\label{sec:geoconsequences}

Table~2 of the main paper compares DreamNav with regressors that use no correspondences, and its margin (83.3\% \SRd{} against at most 11.3\%) does not represent the margin over geometric estimation. Against a classical pipeline assembled from public components and the same flight-height prior, the refined stage gains about five points of success on PairUAV and the coarse stage loses about five points of \SRd{} and fourteen of \SRe; a recent dense matcher is level with the refined stage; and on the real photographs both are more accurate than DreamNav (\cref{sec:real}). The rows of \cref{tab:geometric,tab:real_geo} belong in Tables~2 and~4, and the statements of the main paper about the superiority over existing estimators have to be restricted accordingly (\cref{sec:corrections}). What the experiments support is narrower: pose-guided re-matching at the original resolution improves a learned sparse matcher from the level of \gbSgRotSRd\% to the level of a high-resolution classical pipeline in direction; a learned fusion of several pose and scale estimates is more accurate than any single one of them, including the classical pipeline; and the classical match set can replace the imagined one in this fusion at a small loss (\cref{sec:siftfusion}).
